\documentclass{beamer}

\usepackage{pgf,pgfpages}
\usepackage{graphicx}
\usepackage{units}
\usepackage[utf8]{inputenc}
\usepackage{lmodern}
\usepackage{subfig}
\makeatletter
\newcommand\ChangeItemFont[3]{%
\renewcommand{\itemize}[1][]{%
  \beamer@ifempty{##1}{}{\def\beamer@defaultospec{#1}}%
  \ifnum \@itemdepth >2\relax\@toodeep\else
    \advance\@itemdepth\@ne
    \beamer@computepref\@itemdepth% sets \beameritemnestingprefix
    \usebeamerfont{itemize/enumerate \beameritemnestingprefix body}%
    \usebeamercolor[fg]{itemize/enumerate \beameritemnestingprefix body}%
    \usebeamertemplate{itemize/enumerate \beameritemnestingprefix body begin}%
    \list
      {\usebeamertemplate{itemize \beameritemnestingprefix item}}
      {\def\makelabel####1{%
          {%
            \hss\llap{{%
                \usebeamerfont*{itemize \beameritemnestingprefix item}%
                \usebeamercolor[fg]{itemize \beameritemnestingprefix item}####1}}%
          }%
        }%
  \ifnum\@itemdepth=1\relax
    #1%
  \else
  \ifnum\@itemdepth=2\relax
    #2%
  \else
  \ifnum\@itemdepth=3\relax
    #3%
  \fi%
  \fi%
  \fi%
  }
  \fi%
  \beamer@cramped%
  \raggedright%
  \beamer@firstlineitemizeunskip%
}}
\makeatother

\mode<presentation>
{
  \usetheme{ift}
  \setbeamercovered{transparent}
  \setbeamertemplate{items}[square]
}

\usefonttheme[onlymath]{serif}
\setbeamerfont{frametitle}{size=\LARGE,series=\bfseries}

\definecolor{uibred}{RGB}{102,0, 0}
\definecolor{uibblue}{RGB}{102,0, 0}
\definecolor{uibgreen}{RGB}{102,0, 0}
%\definecolor{uibgreen}{RGB}{50, 105, 0}
\definecolor{uiborange}{RGB}{102,0, 0}


\beamertemplatenavigationsymbolsempty



\defbeamertemplate{enumerate item}{mycircle}
{
  %\usebeamerfont*{item projected}%
  %\usebeamercolor[bg]{item projected}%
  \begin{pgfpicture}{0ex}{0ex}{1.5ex}{0ex}
	%\pgfcircle[fill]{\pgfpoint{0pt}{.75ex}}{1.25ex}
    \pgfbox[center,base]{\color{uibblue}\insertenumlabel.}
  \end{pgfpicture}%
}
[action]
{\setbeamerfont{item projected}{size=\scriptsize}}
\setbeamertemplate{enumerate item}[mycircle]




\title{Assessing Fox Mortality in Fox Penning Trials}
\author{Ian Crandell \& Andy Hoegh}
\institute{
Department of Statistics, Virginia Tech
}
\date{12.4.2014}

\begin{document}

\newcommand\Fontvi{\fontsize{6}{7.2}\selectfont}

\setbeamertemplate{background}
 {\includegraphics[width=\paperwidth,height=\paperheight]{frontpage}}
\setbeamertemplate{footline}[default]

\begin{frame}
  \titlepage
  \vspace{5cm}
\end{frame}

%
% Set the background for the rest of the slides.
% Insert infoline at the end
%
\setbeamertemplate{background}
 {\includegraphics[width=\paperwidth,height=\paperheight]{slide_small}}
\setbeamertemplate{footline}[ifttheme]

%%%%%%%%%%%%%%%%%%%%%%%%%%%%%%%%%%%%%%%%%%%%%%%%%%%%%%
%%%%%%%%%%%%%%%%%%%%%%%%%%%%%%%%%%%%%%%%%%%%%%%%%%%%%%
\section{Motivation}
\subsection{Intro}
\begin{frame}
	\frametitle{What is fox penning?}
	\begin{columns}
		\column{0.5\textwidth}
		\parbox[c][0.9\textheight]{0.9\textwidth}
		{\includegraphics<1>[width=\textwidth]{foxhunting.jpg}}
		\column{0.5\textwidth}
		\parbox[c][0.9\textheight]{0.9\textwidth}
		{\begin{itemize}
				\item Fox hunting is an ancient sport where human trainers and dogs hunt wild foxes.
				\item For fox penning, the foxes are brought into an enclosure and are then hunted for training or entertainment.
				\item Fox penning does not adhere to the principle of \emph{fair chase}.
			\end{itemize}}
	\end{columns}
%	\frametitle{Goal: Predict Civil Unrest}
%	\begin{columns}
%		\column{0.5\textwidth}
%			\parbox[c][0.9\textheight]{0.9\textwidth}
%				{\includegraphics<1>[width=\textwidth]{Venezuela.png}}
%		\column{0.5\textwidth}
%			\parbox[c][0.9\textheight]{0.9\textwidth}
%			{\includegraphics<1>[width=\textwidth]{BrazilProtestjpg.jpg}}
%	\end{columns}
\end{frame}
%%%%%%%%%%%%%%%%%%%%%%%%%%%%%%%%%%%%%%%%%%%%%%%%%%%%%%
\begin{frame}
	\frametitle{Project Relevance}
	\center
			{\includegraphics<1>[width=.8\textwidth]{Media.png}}		
\end{frame}
%%%%%%%%%%%%%%%%%%%%%%%%%%%%%%%%%%%%%%%%%%%%%%%%%%%%%%
\begin{frame}
	\frametitle{Origins of the Project}
	\begin{itemize}
		\item In early 2013 the Virginia Department of Game and Inland Fisheries (VDGIF) was considering new regulations on fox penning. They brought this dataset to Virginia Tech's Conservation Management Institute (CMI). 
		\item The focus was on assessing factors contributing to fox mortality in fox-penning enclosures.
	\end{itemize}
\end{frame}
%%%%%%%%%%%%%%%%%%%%%%%%%%%%%%%%%%%%%%%%%%%%%%%%%%%%%%
%%%%%%%%%%%%%%%%%%%%%%%%%%%%%%%%%%%%%%%%%%%%%%%%%%%%%%
\section{Data Overview}
\subsection{Overview}
\begin{frame}
	\frametitle{Data Overview}
		\begin{itemize}
			\item Data consists of two types of days: trial days and non-trial days from a single fox pen in Virginia. For each day the following information is available:
			\begin{itemize}
				\item Number of dogs (for trial days only)
				\item Experience of each fox 
				\item Survival of each fox
			\end{itemize}
			\item For non-trial days, we knew when a fox died, but only had vague information about the number of dogs. We handled this limitation by taking a missing data approach to the analysis.
		\end{itemize}
	\end{frame}
%%%%%%%%%%%%%%%%%%%%%%%%%%%%%%%%%%%%%%%%%%%%%%%%%%%%%
\subsection{Missing Data}
\begin{frame}
	\frametitle{Missing Data}
	Inferences are made from $P(\Theta|X,Y),$ where $
    \mathcal{X}=\left[
                \begin{array}{ll}
                  X \\
                  X^{*} 
                \end{array}
              \right]$ \\and $X^{*}$ denotes missing covariates. Then in a Bayesian paradigm
              \begin{eqnarray*}
              P(\Theta|X,Y), = \int P(\Theta|X,X^{*},Y) p(X^{*}) dX^{*}
              \end{eqnarray*}
	\end{frame}
%%%%%%%%%%%%%%%%%%%%%%%%%%%%%%%%%%%%%%%%%%%%%%%%%%%%%%
%%%%%%%%%%%%%%%%%%%%%%%%%%%%%%%%%%%%%%%%%%%%%%%%%%%%%%
\section{Modeling}
\subsection{Modeling Specification}
\begin{frame}
	\frametitle{Modeling Specification}
\begin{eqnarray}
y_{it} &\sim& Bernoulli(p_{it}) \\
probit(p_{it}) & = & \alpha + \mathcal{X}_{it}\beta+ \theta_{it} \\
\theta_{it} &\sim& N(0, \phi^{-1})
\end{eqnarray}
where $\mathcal{X}_{it} = [Num.dogs_t\; log.experience_{it}\; (num.dogs*log.experience)_{it}],$ $i = \{1,...27\} $(fox), and $t= \{1,...,T\}$ (trials). Let $R_{it}$ be the risk matrix, where
\[
    R_{it}=\left\{
                \begin{array}{ll}
                  1 \quad \text{ if fox i is alive and collared on day t-1}\\
                  0 \quad \text{otherwise}
                \end{array}
              \right.
  \]
 \end{frame}
%%%%%%%%%%%%%%%%%%%%%%%%%%%%%%%%%%%%%%%%%%%%%%%%%%%%%%
\begin{frame}
	\frametitle{Modeling Specification}
	 We use data augmentation where 
	 \[
    Y_{it}=\left\{
                \begin{array}{ll}
                  1 \quad  Z_{it} > 0\\
                  0 \quad Z_{it} \leq 0
                \end{array}
              \right.
  \]
and then
  \begin{eqnarray}
  Z_{it} \sim N( \alpha + \mathcal{X}_{it}\beta +\theta_{it},1).
  \end{eqnarray}
\\
where vague normal priors are placed on $\alpha, \beta$. Then $p(\phi) \sim Gamma$ and $X^{*} \sim N(0,10) \mathcal{I}(X^* > 0 \;\&\; X^* < 50)$ 
\end{frame}
%%%%%%%%%%%%%%%%%%%%%%%%%%%%%%%%%%%%%%%%%%%%%%%%%%%%%%
%%%%%%%%%%%%%%%%%%%%%%%%%%%%%%%%%%%%%%%%%%%%%%%%%%%%%%
\section{Results}
\subsection{Coefficient Estimates}
\begin{frame}
	\frametitle{Posterior Summaries}
\begin{table}[h]
%\caption{default}
\begin{center}
\begin{tabular}{|c|c|c|}
\hline
 & mean & CI \\
 \hline
  $\alpha$ & 5.4 & (3.0,11.6) \\
 $\beta_{dogs}$ & -.008 & (-.015, -.004) \\
 $\beta_{exp}$ & -.65 & (-2.04,-.10) \\
 $\beta_{dogs*exp}$ & .0014 & (.0006, .0030)\\
 \hline
\end{tabular}
\end{center}
\label{default}
\end{table}%

\end{frame}
%%%%%%%%%%%%%%%%%%%%%%%%%%%%%%%%%%%%%%%%%%%%%%%%%%%%%%
\subsection{Graphical Summaries}
\begin{frame}
	\frametitle{Response Surface}
	\begin{columns}
		\column{0.5\textwidth}
			\parbox[c][0.9\textheight]{0.9\textwidth}
				{\includegraphics<1>[width=\textwidth]{ResponseSurface.pdf}}
		\column{0.5\textwidth}
			\parbox[c][0.9\textheight]{0.9\textwidth}
			{\includegraphics<1>[width=\textwidth]{SurvivalProb.pdf}}
	\end{columns}

\end{frame}
%%%%%%%%%%%%%%%%%%%%%%%%%%%%%%%%%%%%%%%%%%%%%%%%%%%%%%
\begin{frame}
	\frametitle{Regime Study}
	All regimes assume a fox was placed in the pen on the first day of the study.
	\begin{itemize}
	\item Regime A: No constraints on number of dogs or allotted acclimation time.
	\item Regime B: Two weeks acclimation time with no dogs.
	\item Regime C: No more than 400 dogs allowed in pen at a time.
	\item Regime D: Two weeks acclimation time and 400 dog limit.
	\end{itemize}
%	At the finest level $p(\boldsymbol{Y}_{tL_j}|\boldsymbol{\mu}_{tL_j}) \overset{ind}{\sim} \text{Poisson}
\end{frame}


\begin{frame}
	\frametitle{Regime Survival Plots}
	\begin{columns}
		\column{0.5\textwidth}
			\parbox[c][0.9\textheight]{0.9\textwidth}
				{\includegraphics<1>[width=\textwidth]{RegA.pdf}}
						\column{0.5\textwidth}
			\parbox[c][0.9\textheight]{0.9\textwidth}
			{\includegraphics<1>[width=\textwidth]{RegB.pdf}}
	\end{columns}
\end{frame}

\begin{frame}
	\frametitle{Regime Survival Plots}
	\begin{columns}
		\column{0.5\textwidth}
			\parbox[c][0.9\textheight]{0.9\textwidth}
				{\includegraphics<1>[width=\textwidth]{RegC.pdf}}
						\column{0.5\textwidth}
			\parbox[c][0.9\textheight]{0.9\textwidth}
			{\includegraphics<1>[width=\textwidth]{RegD.pdf}}
	\end{columns}
\end{frame}

\begin{frame}
	\frametitle{Regime Analysis}
		\begin{itemize}
			\item Mortality is greatest with many dogs and for inexperienced foxes, but survival greatly improves either with a capped dog count or some time for the foxes to acclimate to the pen.
			\item Trials are run weekly, so two weeks acclimation is impractical. Thus, we can most practically smooth out fox mortality by capping the number of dogs in the pen.
		\end{itemize}
\end{frame}

\subsection{Final Thoughts}
\begin{frame}
	\frametitle{What's Ahead}
		While we have addressed fox mortality, it remains to make the connection between mortality and cruelty. 
\end{frame}

\end{document}

